\section{Colored norm support and a density--window tradeoff}
\label{ncs-section}

For an integer \(B\geq1\), retain the actual block
\[
 I_B=\{k\in\mathbb Z:B\leq k<2B\},\qquad
 a_k=3k,\quad w_k=a_k+\zeta,\quad
 Q_k=a_k^2+a_k+1,\quad \alpha_k=w_k/\bar w_k.
\]
Here \(\zeta^2=\zeta-1\). The roots are primitive and unramified,
and \(9B^2\leq Q_k<49B^2\). For real \(x\geq3\), define
\[
 \mathcal S(B,x)=\{k\in I_B:\text{some prime }p>x\text{ divides }Q_k\}.
\]
This is a condition on the actual norms, fixed before a boundary
prime, depth or tuple length is chosen. All prime sums below range
over rational primes.

\begin{theorem}[An explicit finite independent coloring]\label{ncs-colors}
For every integer \(B\geq1\) and real \(x\geq3\), put
\[
 m=\left\lfloor\frac{\log(49B^2)}{\log x}\right\rfloor,
 \qquad D=1+m\left(2\left\lfloor\frac Bx\right\rfloor+1\right).
\]
Then \(\mathcal S(B,x)\) is a disjoint union of at most \(D\)
classes on each of which the actual ratios \(\alpha_k\) are
multiplicatively independent. The empty case is included.
\end{theorem}
\begin{proof}
A prime dividing \(Q_k\) is neither 2 nor 3. Since
\(a_k^3\equiv1\pmod p\) and \(a_k\not\equiv1\pmod p\), it is
1 modulo 3. The polynomial \(9X^2+3X+1\) has exactly two simple
roots modulo such a prime: its discriminant is \(-27\).

Join distinct indices in \(\mathcal S(B,x)\) if their norms share
some prime above \(x\). A norm has at most \(m\) distinct such
primes, since their product is at most \(Q_k<49B^2\).
For each one of them, the two residue classes meet the block in
at most \(2\lfloor(B-1)/p\rfloor+2\) indices. Removing the given
vertex leaves at most
\[
 2\left\lfloor\frac{B-1}p\right\rfloor+1
 \leq2\left\lfloor\frac Bx\right\rfloor+1
\]
possible neighbors. Thus every vertex has degree at most \(D-1\).
Color vertices in increasing order, always taking the first of
\(D\) colors not used by earlier neighbors. At most \(D-1\)
colors are forbidden, so this finite algorithm succeeds. It does
not invoke a matroid-partition theorem.

Within a color, choose for each root one prime \(p>x\) dividing
its norm. It divides no other norm of that color. At one oriented
prime ideal above \(p\), the ratio \(\alpha_k\) has nonzero
integer valuation \(d_k\). Both conjugate factors cannot occur,
since a common ideal divisor would divide
\(w_k-\bar w_k=\sqrt{-3}\). Every other ratio in the same color
has valuation zero at this ideal. Applying the valuation to a
multiplicative relation gives \(c_kd_k=0\), hence \(c_k=0\),
for each index. The witness need not have depth one; nonzero depth
suffices. Small norm primes may be shared. If \(m=0\), the
norm inequality forces the original vertex set to be empty.
\end{proof}

The colors depend on the actual norm factorizations, \(B,x\),
and the chosen finite ordering, but not on any later boundary
prime. No independent coloring of the norm-smooth complement has
been proved by this argument.

\begin{theorem}[Full valuation mass with its cutoff endpoint]
\label{ncs-small-mass}
There is an absolute constant \(C\) such that, for all integers
\(B\geq3\) and reals \(3\leq x\leq B\), with
\(\Lambda_B=\log(49B^2)\),
\begin{equation}\label{ncs-endpoint}
 \sum_{k\in I_B}\sum_{p\leq x}v_p(Q_k)\log p
 \leq B\log x+CB\log\log(3x)
                    +\frac{Cx\Lambda_B}{\log x}.
\end{equation}
Consequently
\begin{equation}\label{ncs-domain-finite}
 |\mathcal S(B,x)|\geq
 \max\left\{0,
 \frac{2B\log B-B\log x-CB\log\log(3x)-Cx\Lambda_B/\log x}
      {\Lambda_B}\right\}.
\end{equation}
\end{theorem}
\begin{proof}
For a supported prime \(p\), each of the two simple residue roots
lifts uniquely at every depth. Put
\(H_p=\lfloor\Lambda_B/\log p\rfloor\).
Counting the two residue classes at every level gives
\[
 \sum_{k\in I_B}v_p(Q_k)\log p
 \leq \frac{2B\log p}{p-1}+2H_p\log p.
\]
Only \(p\equiv1\pmod3\) occur. The fixed-modulus estimate
\(\vartheta(x;3,1)=x/2+O(x/\log x)\), from Theorem~1.2,
printed page~5, formula~(1.12), of
\href{https://arxiv.org/abs/1802.00085v3}
{Bennett, Martin, O'Bryant and Rechnitzer}, implies by partial
summation
\[
 \sum_{\substack{p\leq x\\p\equiv1\ (3)}}\frac{\log p}{p-1}
       =\frac12\log x+O(\log\log(3x)).
\]
The error integral is \(O(\int^x dt/(t\log t))\), and replacing
\(1/p\) by \(1/(p-1)\) costs a convergent series. Increasing an
absolute constant covers all \(x\geq3\). For the complete endpoint,
\[
 \sum_{\substack{p\leq x\\p\equiv1\ (3)}}2H_p\log p
 \leq2\Lambda_B\pi(x)\ll\frac{x\Lambda_B}{\log x}.
\]
This proves \eqref{ncs-endpoint}. The total actual norm mass is
at least \(2B\log B\); the mass above \(x\) is supported on
\(\mathcal S(B,x)\) and is at most \(\Lambda_B\) per member.
Subtraction and division prove \eqref{ncs-domain-finite}.
\end{proof}

\begin{corollary}[Uniformity on a fixed range of support exponents]
\label{ncs-theta}
Fix \(0<\theta_0\leq1\). Uniformly for
\(\theta\in[\theta_0,1]\), as integer \(B\) tends to infinity,
\begin{align}
 |\mathcal S(B,B^\theta)|
 &\geq(1-\theta/2)B
          -O_{\theta_0}(B\log\log B/\log B),\label{ncs-density}\\
 D&\leq(1+9/\theta_0)B^{1-\theta}.\label{ncs-D}
\end{align}
The second bound holds already for \(B\geq49\) and \(B^\theta\geq3\).
\end{corollary}
\begin{proof}
Here \(\Lambda_B\leq3\log B\), so the endpoint in
\eqref{ncs-endpoint} is at most \(3CB^\theta/\theta_0\).
The main term is \(\theta B\log B\); division by
\(\Lambda_B=2\log B+O(1)\) gives \eqref{ncs-density}.
Also \(m\leq3/\theta_0\) and
\(2\lfloor B^{1-\theta}\rfloor+1\leq3B^{1-\theta}\), proving
\eqref{ncs-D}. No uniform limit as \(\theta_0\) tends to zero
is asserted. This uses the whole block, not an iteration of its
mass estimate on an arbitrary discarded subset.
\end{proof}

For the boundary estimates, let \(n>324\) be prime and \(B=n^4\).
Writing \(w_k^n=A+C\zeta\), put
\[
 T_n(k)=|AC(A+C)|,\qquad
 c_n(k)=\max\{|A|,|C|,|A+C|\},\qquad t_k=\log c_n(k).
\]
A positive-sector unit rotation preserves \(c_n,T_n\).
The elementary actual bounds are
\[
 t_k\geq n\log(3B),\qquad
 \log T_n(k)\leq3n\log(6B+1).
\]
There are no zero boundary coordinates for these primitive
unramified powers.

\begin{theorem}[Uniform moments summed over the actual colors]
\label{ncs-eighty}
For every prime \(n>324\), \(B=n^4\), real \(x\geq3\) and
\(8B<Z\leq B^2\), define
\[
 F_{8B,Z}(k)=\sum_{8B<q\leq Z}(v_q(T_n(k))-3)_+\log q.
\]
With the exact color budget \(D\) from Theorem~\ref{ncs-colors},
\begin{equation}\label{ncs-mean-finite}
 \frac1B\sum_{k\in\mathcal S(B,x)}\frac{F_{8B,Z}(k)}{t_k}
 \leq\frac{80DZ}{Bn\log(Z/(3n))}.
\end{equation}
\end{theorem}
\begin{proof}
We spell out why the uniform-moment proof applies to each color.
For every positive integer \(\nu\), two actual \(\nu\)-fold
products of block roots have coordinates satisfying
\[
 |R|,|R'|\leq2(7B)^\nu,\qquad
 |S|,|S'|\leq2\nu(7B)^{\nu-1}.
\]
The second estimate is the finite telescoping bound for the
imaginary coordinate. For a prime \(q>8B\), the determinant
height divided by \(q^{2\nu}\) is strictly less than
\[
 \frac8{7B}\nu(49/64)^\nu
 <\frac{512}{105B}<1.
\]
The finite geometric sum proves this uniformly for every finite
\(\nu\). Thus congruent ratio products with unit denominators
are exactly equal over \(\mathbb Q(\zeta)\).

Every actual \(q^{2\nu}\)-hit has unit norm and gives a norm-one
\(3n\)-torsion element modulo \(q^{2\nu}\). Since \(q>n\), the
split/inert simple-lifting calculation of Lemma~\ref{mc-torsion}
gives at most \(3n\) such elements at every depth. Within a color,
Theorem~\ref{ncs-colors} recovers the multiplicities from equal
actual products. Thus multisets of \(\nu\) deep roots inject into
this group, giving
\[
 \binom{M_\nu+\nu-1}{\nu}\leq3n
\]
when \(M_\nu>0\), with zero positive-length multisets when it is
empty. In particular \(M_2\leq\sqrt{6n}\). For
\(r=\lceil\sqrt n\rceil\), four roots would give
\(3n\geq\binom{r+3}3>n^{3/2}/6>3n\). Hence \(M_r\leq3\).

Every positive layer is paid by
\((e-3)_+=\sum_{j\geq4}{\bf1}_{e\geq j}\).
In one color the layers \(4\leq j<2r\) contribute at most
\(4\sqrt6\,n\log q<10n\log q\). The remaining layers have at
most three roots and the full height cap gives at most \(9nL\),
where \(L=\log(6B+1)\). Summing colors therefore gives
\[
 \sum_{k\in\mathcal S(B,x)}(v_q(T_n(k))-3)_+\log q
                  \leq D(10n\log q+9nL).
\]
For \(q\leq Z\leq B^2\), both \(\log q\) and \(L\) are at
most \(2\log B\). The full-block normalized cost is at most
\(38D/B\) per prime.

Positive excess has homogeneous rank \(n\). Rank one would,
since \(q>n\), force \(q^2\) into one of the coprime old factors
\(a_k,a_k+1<6B+1\), which is impossible. Thus \(n\) divides
the norm-one group order \(q-\chi(q)\), where
\(\chi(q)=(-3/q)\). That order is also divisible by 3, and
\((n,3)=1\), so \(3n\mid q-\chi(q)\). The order of
\(\alpha_k\) itself may be \(n\) or \(3n\); the argument does
not require the latter. In particular \(q\equiv\pm1\pmod{3n}\).

Montgomery--Vaughan's Brun--Titchmarsh inequality, Theorem~2,
printed page~121, formula~(1.10), in the
\href{https://personal.science.psu.edu/rcv4/personal/Publications/large_sieve.pdf}
{author-hosted primary paper}, bounds these two classes through
\(Z\) by
\[
 \frac{2Z}{(n-1)\log(Z/(3n))}.
\]
Its interval form gives the through-\(Z\) version by decreasing
the left endpoint to zero. Multiply by \(38D/B\), and use
\(76n/(n-1)<80\), to obtain \eqref{ncs-mean-finite}.
\end{proof}

\begin{theorem}[A density--window tradeoff]\label{ncs-tradeoff}
Fix \(3/4<\theta\leq1\) and \(\delta>0\). For sufficiently
large prime \(n\), put \(B=n^4\), \(\kappa=4\theta-3\), and
\[
 Z_n=\left\lfloor Bn^\kappa(\log n)^{1-\delta}\right\rfloor.
\]
Then \(8B<Z_n\leq B^2\), the complete positive-excess mean on
\(\mathcal S(B,B^\theta)\) is \(O_{\theta,\delta}((\log n)^{-\delta})\),
and there is an actual subset \(\mathcal P_{n,\theta,\delta}\)
of \(I_B\) with relative size at least
\[
 1-\frac\theta2-
 O_{\theta,\delta}\!\left(\frac{\log\log n}{\log n}
              +(\log n)^{-\delta/2}+(\log n)^{-1/2}\right)
\]
on which
\[
 k\in\mathcal S(B,B^\theta),\qquad
 \frac{F_{8B,Z_n}(k)}{t_k}\leq(\log n)^{-\delta/2},\qquad
 \frac{L_{8B}(T_n(k))}{t_k}\leq(\log n)^{-1/2}.
\]
Here \(L_X(T)=\sum_{q\leq X}v_q(T)\log q\) is the full mass.
\end{theorem}
\begin{proof}
For fixed \(\kappa>0\), the factor
\(n^\kappa(\log n)^{1-\delta}\) tends to infinity and is
\(o(n^4)\). Thus the cutoff satisfies the asserted domain, and
\[
 \log(Z_n/(3n))=(3+\kappa)\log n
                 +(1-\delta)\log\log n-\log3+o(1).
\]
By \eqref{ncs-D}, \(D=O_\theta(n^{4(1-\theta)})\), and the
power of \(n\) in \eqref{ncs-mean-finite} cancels exactly:
\(4(1-\theta)+\kappa-1=0\). This proves the mean bound.
Markov at \((\log n)^{-\delta/2}\) removes the corresponding
fraction of the entire block. Theorem~\ref{fml-mean}, with
Section~\ref{ebi-section}'s elementary inputs, gives at cutoff
\(8B\) a full-block mean \(O(1/\log n)\); Markov removes
\(O((\log n)^{-1/2})\) more. Intersect these two actual good
sets with \eqref{ncs-density}. All means remain normalized by
\(B\). In particular \((B,8B]\) is paid as full mass, without
a cap imposed there. For fixed \(\theta_0>3/4\) the constants
and sufficiently-large threshold are uniform in
\(\theta\in[\theta_0,1]\), with fixed \(\delta\).
\end{proof}

Define \(J(T)=\log T-3\log\operatorname{rad}(T)\) and let
\(W_Z(T)\) be the signed contribution of supported primes above
\(Z\). On the preceding set, with
\(\varepsilon_n=(\log n)^{-\delta/2}+(\log n)^{-1/2}\),
exact disjoint prime decomposition gives
\begin{align}
 \frac{J(T_n(k))}{t_k}
 &\leq\frac{W_{Z_n}(T_n(k))}{t_k}+\varepsilon_n,
       \label{ncs-signed}\\
 \frac{\log\operatorname{rad}(T_n(k))}{t_k}
 &\geq1-\frac4{3n}-\frac{\varepsilon_n}{3}
                 -\frac{(W_{Z_n}(T_n(k)))_+}{3t_k}.
       \label{ncs-radical}
\end{align}
The first inequality uses the full low mass and the middle positive
excess. The second uses the elementary block bound
\(\log T_n(k)/t_k\geq3-4/n\). The far signed term keeps its
negative credit; its positive part remains an unproved membership
gate, not a consequence of this finite-window average.

The exponent \(\kappa\in(0,1]\) has proved domain fraction
\((5-\kappa)/8\). For example \(\theta=15/16\) gives
\(\kappa=3/4\), fraction \(17/32\), and cutoff
\(\lfloor Bn^{3/4}(\log n)^{1-\delta}\rfloor\).
For \(\kappa<1/2\), the earlier MC result already gives an
intermediate-window domain of fraction tending to one, so this
coloring is additional structure rather than improved coverage
there. The useful coverage tradeoff is \(1/2<\kappa<1\), beyond
the MC power range and with a fraction larger than one half.
No unproved endpoint comparison at \(\kappa=1/2\) is asserted.

The norm-smooth complement, its possible rank deficiency, and the
unbounded farther signed tail remain open. Exact pair-product
collisions disprove a universal one-color assertion, but do not
refute a stronger multi-color theorem for all actual roots.
The full-block norm-mass proof cannot simply be iterated on an
arbitrary discarded subset. The coloring and ideal witnesses above
are ordinary proofs; no new Lean formalization of them or of the
analytic domain theorem is claimed by this section.
